\scanpage{179}\mspage{34}%
% archive sheet number: 179
which are manifolds with boundary, the \deleted{\unsure{morphisms|morphisms}} arrows of the
diagrams being for the most part
locally trivial fibrations of manifolds with boundary over others
(with fibers manifolds with boundary), \uline{and} \deleted{\unsure{some|the}}\unsure{the|some} inclusions
of the boundaries into the manifolds with boundary \deleted{(}[in
fact, one finds slightly
more general inclusions, some of the boundaries which appear
being equipped with an
``elementary'' cellular decomposition, \deleted{into \unsure{cells|elements}}\inserted{i.e. of which \unsure{all}} the
% L9: the live replacement runs above the cancelled baseline; toutes is raised with it.
\inserted{closed} strata are again manifolds with boundary which
are glued together along common parts of the
boundary~; and there is reason to consider \inserted{also} the inclusions
of these pieces into one
another\ldots.\deleted{)}], and is reconstituted from
this diagram \inserted{by gluing}. In other words, one
has a canonical description-dévissage
of compact tame spaces \inserted{$X$}, possibly
equipped with equisingular stratifications
with nonsingular strata, in terms of \inserted{\unsure{finite}} diagrams
of a \unsure{specified|specified} nature formed with manifolds
with boundary. When one is interested in $\unsure{\mathcal{P}|P}$-constructible sheaves
\inserted{(or complexes of sheaves, or $n$-stacks\ldots)}
on $X$, where $\mathcal{P}$ is such a
fixed stratification, these are made explicit in
terms of the diagram under consideration, of which it
